Для сравнения скорости работы аллокаторов была написана несложная тестовая функция, которая создаёт массив из $n$ элементов ($n$ подаётся на вход функции), являющихся указателями на числа типа \texttt{int}. С помощью метода \texttt{allocate}, который реализует поданный на вход функции аллокатор, производится выделение памяти под каждый из объектов, а затем ему присваивается некоторое тривиально вычисляемое значение. После этого производится очистка памяти в определённом порядке: для стекового аллокатора это порядок, обратный выделению, а для пуллового был сделан выбор в пользу случайного порядка освобождения, ведь возможность освобождения произвольного блока памяти "--- одно из главных преимуществ данного аллокатора. Дополнительно вызывается метод \texttt{reset}, который в случае линейного аллокатора полностью его сбрасывает.

\small
\begin{minted}{cpp}
enum AllocType { Linear, Stack, Pool, Standard };

template<typename Allocator>
void test(Allocator& allocator, int n,
		  chrono::duration<double> &overhead, AllocType type) {
    int* array[n];
    for (int i = 0; i < n; i++) {
        array[i] = (int*)allocator.allocate(sizeof(int));
        if (array[i]) *array[i] = 6 * i + 7;
    }
	auto start = chrono::high_resolution_clock::now();
  	vector<int> order(n);
	iota(order.begin(), order.end(), 0);
	// для пуллового и стандартного аллокатора
	// тестируем освобождение памяти в произвольном порядке
	if (type == Pool || type == Standard)
		shuffle(order.begin(), order.end(), rnd);
	// для стекового аллокатора тестируем
	// освобождение памяти в порядке, обратном выделению
	else if (type == Stack)
		reverse(order.begin(), order.end()); 
	auto end = chrono::high_resolution_clock::now();
	// лишние операции не участвуют в измерении времени
	overhead = end - start;
	for (int x : order) {
		allocator.deallocate(array[x]);
	}
	allocator.reset();
	// тестируем, что память была освобождена корректно
	array[10] = (int*)allocator.allocate(sizeof(int));
	if (array[10]) *array[10] = 1024;
}
\end{minted}
\normalsize

Для замера времени использовалась функция \texttt{chrono::high\_resolution\_clock::now()}, которая предоставляет доступ к наиболее точным системным часам \cite{cpp-reference}. Время работы функции было вычислено как разность между временем до и после выполнения программы, причём дополнительно был посчитан \texttt{overhead} "--- время, не связанное с аллокацией и деаллокацией, которое не должно было повлиять на измерение и потому вычиталось из общей продолжительности. Программный код тестирующей функции и функции, которая замеряет время работы программы, представлен в приложении А.

В тестах использовалось $N = 1000000$, для каждого аллокатора было выполнено $5$ замеров, чтобы получить более точный усреднённый результат. В ходе эксперимента были получены следующие результаты, которые представлены в таблице \ref{tab:results} и в диаграмме \ref{fig:results}.

\small
\input{table_1.tex}
\normalsize
\input{diagram_1.tex}

Вполне предсказуемым оказалось то, что линейный и стековый аллокатор продемонстрировали наиболее хороший результат. Их действия при аллокации тривиальны: выполняется сравнение, операция арифметики указателей и обычной арифметики. Однако стековый аллокатор также добавляет выделенный блок памяти в список, поэтому было неожиданно увидеть, что он отработал так же быстро, как и линейный аллокатор. Пулловый аллокатор также использует операции с односвязным списком, помимо этого при освобождении памяти применяется дорогостоящая операции взятия по модулю для проверки корректности адреса начала блока. Однако основное замедление, вероятно, связано с произвольным порядком освобождения памяти.

Дополнительно хотелось бы отметить, что стандартный аллокатор иногда демонстрировал совершенно неожиданные результаты: так, при среднем времени работы в $\approx 0.36s$ в один из запусков на выделение памяти ушло более $3s$. Предположительно, это может быть связано с так называемыми <<malloc spikes>> "--- ситуациями, когда на компьютере запущено большое количество процессов, потребляющих оперативную память, и в результате их длительной работы память оказывается сильно фрагментированной. Поиск свободных блоков памяти и объединение маленьких фрагментов памяти в большие связанные блоки (что как раз и делает \texttt{new}/\texttt{malloc}) может вызвать неоправданную задержку при выполнении аллокации. В играх это может привести к внезапному кратковременному зависанию игры, в высоконагруженных приложениях "--- к резкому падению производительности. Такие ситуации, несмотря на кажущуюся редкость, могут вызвать недовольство пользователей, а единственным способом их избежать может оказаться разработка собственного аллокатора \cite{pre-allocation}. Таким образом, исследование ещё раз подтверждает, что далеко не всегда стандартные решения языка C++ являются лучшими для выделения памяти.
